\begin{abstract}
This study proposes an explainable AI–based decision support system to prioritize agricultural skill development among smallholder farmers in Thailand. Building on evidence that farmer capability heterogeneity strongly influences technology adoption, resource-use efficiency, and productivity outcomes \citep{feder1985adoption, davis2012impact}, we analyze observational training data from $n = [N]$ farmers collected between [YEAR--YEAR] with pre- and post-program skill assessments. 

K-Means clustering identifies three distinct farmer segments, consistent with prior segmentation studies in agricultural extension. A Random Forest model is trained to predict skill improvement, and SHAP values are computed to uncover non-linear, segment-specific drivers of training effectiveness \citep{adadi2018xai, molnar2022interpretable}. The model explains approximately 13\% of the variance in skill gain ($R^2 = 0.132$). 

By integrating cluster-level SHAP analysis with a surrogate decision tree, the framework translates complex model behavior into transparent, segment-specific decision rules that support more targeted and accountable training allocation within Thailand’s Smart Farmer policy context.
\end{abstract}
